\textbf{Phần 1. Trắc nghiệm lựa chọn}
\setcounterpageref{ex}{0}

\begin{ex} %[2H6V31]
Đồng hồ chỉ $4$ giờ $15$ phút. Kim phút đang chỉ vào số mấy?
\choice
{Số $4$}
{Số $15$}
{\True Số $3$}
{Số $1$}
\loigiai{
Mỗi khoảng giữa hai số là 5 phút. $15$ phút tương ứng kim phút chỉ vào số $3$.
\textbf{Đáp án của câu hỏi là C}
}
\end{ex}

\begin{ex} %[2H6V31]
Nhìn vào tờ lịch tháng $5$, ngày sinh nhật Bác Hồ là ngày $19$ tháng $5$. Nếu ngày $12$ tháng $5$ là thứ Bảy thì sinh nhật Bác là thứ mấy?
\choice
{\True Thứ Bảy}
{Chủ nhật}
{Thứ Sáu}
{Thứ Hai}
\loigiai{
$19 - 12 = 7$ ngày (đúng 1 tuần). Vậy vẫn là thứ Bảy.
\textbf{Đáp án của câu hỏi là A}
}
\end{ex}

\begin{ex} %[2H6V31]
Đồng hồ điện tử hiển thị $18:30$. Điều đó có nghĩa là mấy giờ?
\choice
{$8$ giờ $30$ phút tối}
{$18$ giờ $30$ phút sáng}
{$6$ giờ chiều}
{\True $6$ giờ $30$ phút tối}
\loigiai{
$18:30$ là $6$ giờ $30$ phút tối (hoặc $6$ giờ rưỡi chiều).
\textbf{Đáp án của câu hỏi là D}
}
\end{ex}

\begin{ex} %[2H6V31]
Lan đi bộ từ nhà đến trường hết $20$ phút. Lan đến trường lúc $7$ giờ $20$ phút sáng. Hỏi Lan xuất phát từ nhà lúc mấy giờ?
\choice
{$6$ giờ $50$ phút}
{\True $7$ giờ sáng}
{$7$ giờ $10$ phút}
{$7$ giờ}
\loigiai{
$7$ giờ $20$ phút $- 20$ phút $= 7$ giờ sáng.
\textbf{Đáp án của câu hỏi là B}
}
\end{ex}

\begin{ex} %[2H6V31]
Khi kim giờ và kim phút của đồng hồ cùng chỉ vào số $12$, lúc đó là mấy giờ?
\choice
{$12$ giờ $30$ phút}
{$6$ giờ}
{\True $12$ giờ đúng}
{$1$ giờ}
\loigiai{
Cả hai kim cùng chỉ số 12 là đúng 12 giờ.
\textbf{Đáp án của câu hỏi là C}
}
\end{ex}

\begin{ex} %[2H6V31]
Trong một tháng, ngày Chủ nhật đầu tiên là ngày $2$. Hỏi ngày Chủ nhật cuối cùng của tháng đó là ngày mấy nếu tháng đó có $30$ ngày?
\choice
{\True Ngày $30$}
{Ngày $28$}
{Ngày $23$}
{Ngày $29$}
\loigiai{
Các ngày Chủ nhật: $2, 9, 16, 23, 30$. Ngày cuối cùng là ngày $30$.
\textbf{Đáp án của câu hỏi là A}
}
\end{ex}

\begin{ex} %[2H6V31]
Mẹ đi chợ từ $8$ giờ sáng đến $9$ giờ $30$ phút sáng. Mẹ đi chợ trong bao lâu?
\choice
{$1$ giờ}
{$2$ giờ}
{$45$ phút}
{\True $1$ giờ $30$ phút}
\loigiai{
$9$ giờ $30$ phút $- 8$ giờ $= 1$ giờ $30$ phút.
\textbf{Đáp án của câu hỏi là D}
}
\end{ex}

\begin{ex} %[2H6V31]
Đồng hồ đang chỉ $2$ giờ rưỡi chiều. Sau $30$ phút nữa thì đồng hồ chỉ mấy giờ?
\choice
{$2$ giờ chiều}
{\True $3$ giờ chiều}
{$3$ giờ $30$ phút}
{$4$ giờ chiều}
\loigiai{
$2$ giờ rưỡi chiều là $2$ giờ $30$ phút chiều. Thêm $30$ phút nữa là đúng $3$ giờ chiều.
\textbf{Đáp án của câu hỏi là B}
}
\end{ex}

\begin{ex} %[2H6V31]
Nếu ngày $28$ tháng $4$ là thứ Bảy thì ngày $1$ tháng $5$ cùng năm đó là thứ mấy? (Biết tháng $4$ có $30$ ngày)
\choice
{Thứ Hai}
{Chủ nhật}
{\True Thứ Ba}
{Thứ Tư}
\loigiai{
Ngày 28: thứ Bảy; Ngày 29: Chủ nhật; Ngày 30: thứ Hai; Ngày 1 tháng 5: thứ Ba.
\textbf{Đáp án của câu hỏi là C}
}
\end{ex}

\begin{ex} %[2H6V31]
Một tiết mục xiếc kéo dài $45$ phút, kết thúc lúc $10$ giờ sáng. Tiết mục xiếc đó bắt đầu lúc mấy giờ?
\choice
{\True $9$ giờ $15$ phút sáng}
{$9$ giờ}
{$9$ giờ $30$ phút}
{$8$ giờ $45$ phút}
\loigiai{
$10$ giờ trừ đi $45$ phút là $9$ giờ $15$ phút sáng.
\textbf{Đáp án của câu hỏi là A}
}
\end{ex}

\begin{ex} %[2H6V31]
Đồng hồ có kim phút chỉ số $12$, kim giờ chỉ số $7$. Lúc này là
\choice
{$12$ giờ $7$ phút}
{$7$ giờ $12$ phút}
{$6$ giờ $50$ phút}
{\True $7$ giờ đúng}
\loigiai{
Kim phút chỉ số 12, kim giờ chỉ số 7 là đúng 7 giờ.
\textbf{Đáp án của câu hỏi là D}
}
\end{ex}

\begin{ex} %[2H6V31]
Nếu hôm nay là thứ Sáu ngày $13$. Thứ Sáu tuần trước là ngày mấy?
\choice
{Ngày $7$}
{\True Ngày $6$}
{Ngày $5$}
{Ngày $8$}
\loigiai{
$13 - 7 = 6$.
\textbf{Đáp án của câu hỏi là B}
}
\end{ex}

\textbf{Phần 2. Trắc nghiệm đúng sai}
\setcounterpageref{ex}{0}

\begin{ex} %[2H6V31]
Xét tính đúng sai của các tình huống thực hành xem giờ và lịch:
\choiceTF
{\True Khi kim phút chỉ vào số 3 là biểu thị 15 phút.}
{Khi kim phút chỉ vào số 6 là biểu thị 45 phút.}
{\True Hai ngày Chủ nhật liên tiếp trong cùng một tháng luôn cách nhau đúng 7 ngày.}
{\True Từ 8 giờ đến 9 giờ kim phút quay trọn vẹn đúng 1 vòng.}
\loigiai{
\begin{itemize}
    \item Ý a) Đúng: $3 \times 5 = 15$ phút. Đúng.
    \item Ý b) Sai: Kim phút chỉ số 6 là 30 phút.
    \item Ý c) Đúng: Đúng, chu kỳ 1 tuần là 7 ngày.
    \item Ý d) Đúng: 1 giờ ứng với 1 vòng kim phút.
\end{itemize}
\begin{itemize}
    \item a) \textbf{Đúng}.
    \item b) \textbf{Sai}.
    \item c) \textbf{Đúng}.
    \item d) \textbf{Đúng}.
\end{itemize}
}
\end{ex}

\begin{ex} %[2H6V31]
Lịch thi học kì của bạn Minh diễn ra trong tuần:
\choiceTF
{\True Thứ Hai ngày 10 thi Toán, thứ Năm cùng tuần thi Tiếng Việt thì thi Tiếng Việt vào ngày 13.}
{\True Nếu môn Toán thi lúc 8 giờ và làm bài 40 phút thì kết thúc lúc 8 giờ 40 phút.}
{Môn Tiếng Việt thi từ 9 giờ đến 10 giờ 30 phút là thi trong 2 giờ.}
{Thứ Bảy tuần đó là ngày 14.}
\loigiai{
\begin{itemize}
    \item Ý a) Đúng: $10 + 3 = 13$. Đúng.
    \item Ý b) Đúng: Đúng.
    \item Ý c) Sai: Chỉ có 1 giờ 30 phút.
    \item Ý d) Sai: Thứ Hai ngày 10, thứ Bảy là $10 + 5 = 15$ (không phải 14).
\end{itemize}
\begin{itemize}
    \item a) \textbf{Đúng}.
    \item b) \textbf{Đúng}.
    \item c) \textbf{Sai}.
    \item d) \textbf{Sai}.
\end{itemize}
}
\end{ex}

\textbf{Phần 3. Trả lời ngắn}
\setcounterpageref{ex}{0}

\begin{ex} %[2H6V31]
Bác Hùng đi ô tô từ 14 giờ đến 17 giờ thì đến nơi. Bác đi hết bao nhiêu giờ?
\shortans{3}
\loigiai{
\textbf{Đáp số:} $3$\\
$17 - 14 = 3$ giờ.
}
\end{ex}

\begin{ex} %[2H6V31]
Đồng hồ chỉ 9 giờ rưỡi, tức là 9 giờ bao nhiêu phút?
\shortans{30}
\loigiai{
\textbf{Đáp số:} $30$\\
Rưỡi là 30 phút.
}
\end{ex}

\begin{ex} %[2H6V31]
Tháng Ba có 31 ngày. Hỏi tháng Ba có bao nhiêu ngày nhiều hơn một tháng có 30 ngày?
\shortans{1}
\loigiai{
\textbf{Đáp số:} $1$\\
$31 - 30 = 1$ ngày.
}
\end{ex}

\begin{ex} %[2H6V31]
Nếu thứ Tư tuần này là ngày 15 thì thứ Tư tuần trước là ngày mấy?
\shortans{8}
\loigiai{
\textbf{Đáp số:} $8$\\
$15 - 7 = 8$.
}
\end{ex}

\begin{ex} %[2H6V31]
Một trận thi đấu cờ vua bắt đầu lúc 8 giờ 10 phút và kết thúc lúc 8 giờ 40 phút. Trận đấu diễn ra trong bao nhiêu phút?
\shortans{30}
\loigiai{
\textbf{Đáp số:} $30$\\
$40 - 10 = 30$ phút.
}
\end{ex}

\begin{ex} %[2H6V31]
Nếu ngày 25 của một tháng là thứ Bảy thì thứ Bảy tuần trước là ngày mấy?
\shortans{18}
\loigiai{
\textbf{Đáp số:} $18$\\
$25 - 7 = 18$.
}
\end{ex}
